\chapter{Conclusions}\label{ch:conclusion}

This project set out to connect an unpublished \textit{Moringa oleifera} irradiation study to
a defensible computational investigation of cellular radiosensitivity, and to use the result
to identify compounds worth testing.

\paragraph{What was established.} A reproducible, provenance-tracked pipeline links public
clonogenic radiosensitivity data to baseline expression for 237 independent biological
components, with cell-line relatedness handled correctly and all preprocessing confined to
training folds. Within that design, pathway-level expression features reduce the error of a
tissue-only baseline by 0.00383 mSF2 units (2.85\%), consistently in direction across nine of
ten cross-validation repeats but small, non-uniform across tissues and unvalidated by any
formal test. Overall predictive power is weak: out-of-sample $R^2$ is near zero for every
model, closely reproducing the strongest published negative benchmark in the field.

Three further analyses of the same cohort tested how robust that advantage is. A random
forest did worse than the linear offset model; on transformed outcome scales the advantage
kept its direction but nearly vanished on one; and an expression-only model transferred to
unseen tissues with the only clearly positive out-of-sample $R^2$ of the project, but lost to
the global mean when Pancreas or Esophagus/Stomach was held out.

The compound module produced a curated, fully referenced evidence base on \textit{Moringa}
constituents in which the extraction solvent is tracked explicitly, and a shortlist of 18
candidates with 5 in tier~A, in which the prespecified liability cap and the not-assessable
rule both visibly bound. It established that niazimicin is the only constituent confirmed in a
dichloromethane fraction of seed and leaf, and that moringin, the isothiocyanate most
characteristic of the plant, has no published ionising-radiation study at all.

\paragraph{What remains hypothetical.} Everything downstream of the prediction. The pathway
ranking is capped at an exploratory tier because no prediction gain was formally shown and no
negative controls were run. The cell-line discrepancies are review signals dominated by
shrinkage. No compound has been shown to be present in the extracts that were tested, no
compound has been shown in this work to modify clonogenic survival, and no claim of selective
radioprotection is made or supported. The candidate tiers describe the state of the published
evidence, not expected experimental outcomes.

\paragraph{Contribution.} The project's contribution is methodological and evidential rather
than predictive: a transparent framework; an honest, negative-leaning benchmark that agrees
with the best prior work; the correction of its own earlier artefact; an explicit record of a
deadline-driven protocol deviation; and a traceable path from the original radiobiology to a
small set of experiments that would settle the question. Its most valuable single output may
be the mismatch revealed by the compound curation --- the compound present in the right
solvent has no radiation data, and the compounds with radiation data were never found in that
solvent --- because it turns a diffuse phytochemical question into two specific and feasible
experiments: characterise the original extracts, and test the best-documented candidates on
clonogenic survival in matched normal and tumour cells.
